% Tables from tools/make_paper_tables.py (evaluation/results/*.csv, run 20260908_194424).
\subsection{Objective Metrics}
\label{objective_metrics}

The objective evaluation quantifies speaker similarity, intelligibility and
response time without human raters. The patient's own recordings are included
as a fourth condition (REAL). This condition is a reference point, not a system
under test: a cosine similarity has no absolute scale, so a value can only be
judged against what two genuine recordings of the same speaker reach under the
same measurement. The recording that serves as the similarity reference is
excluded from REAL, since its similarity to itself is $1.0$ by construction,
leaving $15$ of $16$ texts.

\textbf{Speaker Similarity (SIM)} is the cosine similarity between ECAPA-TDNN
speaker embeddings \cite{desplanques2020ecapa} of the rendered audio and of a
held-out genuine recording of the patient. This reference is a different voice
message from the one used for enrolment, so a system cannot score by
reproducing its own prompt (SpeechBrain \texttt{spkrec-ecapa-voxceleb}
\cite{ravanelli2021speechbrain}); higher is better
\cite{du2025cosyvoice3inthewildspeech}. The model is trained on English speech
and is out of domain for German phone recordings, which affects all values, not
only their differences.

\textbf{Word Error Rate (WER)} is computed from Whisper \cite{radford2023whisper}
(\texttt{small}) transcriptions against the ground-truth text, after identical
normalisation of both sides (case, punctuation, hyphens, umlauts, digits to
German number words). A high WER indicates unclear articulation
\cite{du2025cosyvoice3inthewildspeech}.

\textbf{Service real-time factor (service RTF)} is the wall-clock latency of the
API call divided by the duration of the returned audio
\cite{achanta2021ondeviceneuralspeechsynthesis}. It is termed \emph{service}
RTF because it measures hosted services over the internet, including network
transfer, and is not comparable with the real-time factors reported by the
developers for local inference.

Each text was synthesised three times per system and analysed as returned by
the service, without resampling or normalisation. Repetitions are averaged per
text; the reported statistic is the median across texts with a percentile
95\,\% confidence interval from $4000$ bootstrap replicates over texts.
Differences between conditions are paired by text and judged at a
Bonferroni-adjusted level, since four conditions give six comparisons per
metric and the family-wise error rate would otherwise reach approximately
$26\,\%$. A difference
whose adjusted interval contains zero is reported as \emph{not detectable},
which is a statement about the resolution of the measurement, not about
equality. Qwen3 and S2 Pro rendered all $16$ texts usably; three Voxtral
renders contain no measurable pause floor and were excluded from the listening
test under the rule fixed in advance (Section~\ref{subjective_evaluation}).

\begin{table}[H]
    \centering
    \small
    \begin{tabular}{lccc}
        \toprule
        \textbf{Condition} & \textbf{SIM} $\uparrow$ & \textbf{WER} $\downarrow$ & \textbf{Service RTF} $\downarrow$ \\
        \midrule
        Qwen3 & 0.755 [0.738, 0.769] & 0.097 [0.075, 0.115] & 0.392 [0.377, 0.410] \\
        S2 Pro & 0.703 [0.687, 0.713] & 0.119 [0.101, 0.142] & 0.227 [0.216, 0.249] \\
        Voxtral & 0.609 [0.602, 0.627] & 0.116 [0.107, 0.141] & 0.143 [0.107, 0.158] \\
        \midrule
        Real & 0.786 [0.760, 0.825] & 0.190 [0.147, 0.228] & --- \\
        \bottomrule
    \end{tabular}
    \caption{Objective results, reported as median [95\,\% CI] across texts; intervals from 4,000 bootstrap replicates over texts. $n = 16$ texts per system, $n = 15$ for the reference condition (Real), which excludes the recording used as the similarity reference. Service RTF applies to the three services only.}
    \label{tab:objective_results}
\end{table}


Table~\ref{tab:objective_results} shows the results. In similarity the systems
rank Qwen3, S2 Pro, Voxtral, and all differences between systems are
detectable. The difference between Qwen3 and the genuine recordings ($-0.031$,
adjusted CI $[-0.068, +0.004]$) is not detectable, whereas S2 Pro and Voxtral
remain measurably below the reference. This describes the resolution of
ECAPA-TDNN on this material, not a faithful reproduction of the patient's
voice.

All three systems reach a lower WER than the genuine recordings ($0.097$ to
$0.119$ against $0.190$). This is not a gain in intelligibility: the systems
read a prepared transcript, whereas REAL is spontaneous phone speech. The
difference between Qwen3 and S2 Pro is not detectable ($-0.029$, adjusted CI
$[-0.061, +0.001]$).

Response time ranks in the reverse order: Voxtral is fastest (service RTF
$0.143$), Qwen3 slowest ($0.392$), with all pairwise differences detectable.
All three stay below $1.0$ and thus generate audio faster than it is spoken.
